\documentclass[border=8pt]{standalone}
\usepackage{circuitikz}
\begin{document}
\begin{circuitikz}[american voltages]
\draw
  (3,5) node[vcc]{$V_{DD}$}
  to [short] (3,4.4)
  node[pmos, anchor=S] (MP1) {$M_{P1}$}
  (MP1.D) to [short] (3,2.75)
  to [short] (3,2.6)
  node[nmos, anchor=D] (MN1) {$M_{N1}$}
  (MN1.S) to [short] (3,0.8)
  node[ground]{};
\draw
  (5.4,5) node[vcc]{$V_{DD}$}
  to [short] (5.4,4.4)
  node[pmos, anchor=S] (MP2) {$M_{P2}$}
  (MP2.D) to [short] (5.4,2.75)
  to [short] (5.4,2.6)
  node[nmos, anchor=D] (MN2) {$M_{N2}$}
  (MN2.S) to [short] (5.4,0.8)
  node[ground]{};
\draw
  (3,2.75) to [short, *-] (4.2,2.75)
  node[above]{$v_{mid}$}
  (4.2,2.75) to [short] (4.2,3.5)
  (MP2.G) to [short] (4.2,3.5)
  (MN2.G) to [short] (4.2,2.0)
  (4.2,3.5) to [short] (4.2,2.0);
\draw
  (MP1.G) to [short] (1.5,3.5)
  (MN1.G) to [short] (1.5,2.0)
  (1.5,3.5) to [short] (1.5,2.0)
  to [short, -o] (0.8,2.75)
  node[left]{$v_{in}$};
\draw
  (5.4,2.75) to [short, *-o] (7.2,2.75) node[right]{$v_{out}$};
\end{circuitikz}
\end{document}
